%!BIB program = bibtex
\documentclass{article}
\usepackage[margin=3cm]{geometry}
\usepackage{amsmath}
\usepackage{amssymb}
\usepackage{physics}

\newcommand{\bb}[1]{\mathbb{#1}}
\renewcommand{\vec}[1]{\boldsymbol{#1}}
\renewcommand{\order}[1]{\mathcal{O}\big(#1\big)}

\title{Symplectic integrator}
\author{}
\date{}
\begin{document}

\maketitle

\section{Hamiltonian mechanics}

\subsection{Generalized Momentum}
The generalized momentum $p_i$ corresponding to the generalized coordinate $q^i$ is given by:
\[ p_i = \pdv{L}{\dot{q}^i} \]
where $L$ is the Lagrangian.

\subsection{Hamiltonian}
Hamiltonian $H$ is given by:
\[ H(q,p,t) = \sum_i p_i \, \dot{q}^i(q,p) - L(q,\dot{q}(q,p),t) \]
where $L(q,\dot{q}(q,p),t)$ is the Lagrangian.

\subsection{Hamilton's Equations}
\begin{align*}
	\dot{q}^i &= \pdv{H}{p_i} \\
	\dot{p}_i &= -\pdv{H}{q^i}
\end{align*}

Or in vector notation,
\[ \vec{z} = 
\begin{pmatrix}
	\vec{q} \\
	\vec{p}	
\end{pmatrix} \, , \
J = 
\begin{pmatrix}
	0 & I \\
	-I & 0
\end{pmatrix} \]

\[ \dot{\vec{z}} = J \grad_{\vec{z}}{H(\vec{z})} \]

\section{Poisson Bracket}

\[ \{ A(q,p,t) , B(q,p,t) \} = \sum_i \left( \pdv{A}{q^i} \pdv{B}{p_i} - \pdv{B}{q^i} \pdv{A}{p_i} \right) \]

\[ \dv{A(q,p,t)}{t} = \{ A,H \} + \pdv{A}{t} \]

\section{Symplectic Condition}

Under a canonical transformation ($\vec{z} \to \vec{w}$) :
\begin{align*}
	\dot{\vec{z}} &= J \grad_{\vec{z}}{H(\vec{z})} \\
	\dot{\vec{w}} &= J \grad_{\vec{w}}{H(\vec{w})}
\end{align*}
The Jacobian matrix $M$ is defined as
\[ M_{ij} = \pdv{w_i}{z_j} \]
Using the chain rule:
\begin{align*}
	(\grad_{\vec{z}}{H})_i &= \pdv{H}{z_i}\\
	\pdv{H}{z_i} &= \sum_j \pdv{H}{w_j} \pdv{w_j}{z_i} \\
	&= \sum_j (M^T)_{ij} (\grad_w H)_j \\
	\grad_{\vec{z}}{H(\vec{z})} &= M^T \grad_w{H(\vec{w})}
\end{align*}

\begin{align*}
	\dv{w_i}{t} &= \sum_j \pdv{w_i}{z_j} \dv{z_j}{t} \\
	\dot{w_i} &= \sum_j M_{ij} \dot{z_j} \\
	\dot{\vec{w}} &= M \dot{\vec{z}} \\
	&= M J \grad_{\vec{z}}{H(\vec{z})} \\
	&= M J M^T \grad_{\vec{w}}{H(\vec{w})}
\end{align*}
Canonical transformation $\iff$ the symplectic condition:
\[ M J M^T = J \]

\section{Splitting methods for separable Hamiltonians}
\[ \hat{D}_f = \{ \cdot , f \} \]
\begin{equation} \label{eq:Hamilton_eq_simplified}
	\dot{\vec{z}} = \{ \vec{z}, H \} = \hat{D}_H \vec{z}
\end{equation}
By Taylor expansion:
\begin{align*}
	\vec{z}(t+\Delta t) &= \sum_{n=0}^{\infty} \frac{1}{n!} \dv[n]{\vec{z}(t)}{t} \left( \Delta t \right)^n \\
	&= \sum_{n=0}^{\infty} \frac{1}{n!}(\Delta t)^n \hat{D}_H^n \, \vec{z}(t) \\
	&= \exp(\Delta t \hat{D}_H) \vec{z}(t)
\end{align*}

Thus, the exact solution of \eqref{eq:Hamilton_eq_simplified} is given by:
\[ \vec{z}(t) = \exp(t \hat{D}_H) \vec{z}(0) \]

If the Hamiltonian is separable, 
\[ H(\vec{q},\vec{p}) = K(\vec{p}) + U(\vec{q}) \]

the operator splits as:
\begin{align*}
	\hat{D}_H &= \{ \cdot , H \} =\{ \cdot , K+U \} \\
	&= \{ \cdot, K\} + \{ \cdot, U\} \\
	&= \hat{D}_K + \hat{D}_U
\end{align*}
\[ \exp(\Delta t \hat{D}_H) = \exp(\Delta t (\hat{D}_K + \hat{D}_U) ) = \exp(\Delta t \hat{D}_K + \Delta t \hat{D}_U ) \]

Because $\hat{D}_K$ and $\hat{D}_U$ are non-commutative operators ($[\hat{D}_K, \hat{D}_U] \neq 0$),
\[ \exp(\Delta t \hat{D}_K + \Delta t \hat{D}_U ) \neq \exp(\Delta t \hat{D}_K) \exp(\Delta t \hat{D}_U) \]
(refer to the Baker--Campbell--Hausdorff formula)

Instead, the time-evolution operator $\exp(\Delta t (\hat{D}_K + \hat{D}_U) )$ can be approximated by a product of operators:
\begin{equation}\label{eq:time-evolution_op_prod_approx}
	\exp(\Delta t (\hat{D}_K + \hat{D}_U) ) = \prod_{i=k,\cdots,1} \exp(d_i \Delta t \hat{D}_K) \exp(c_i \Delta t \hat{D}_U) + \order{(\Delta t)^{k+1}}
\end{equation}
where $k\in \bb{N}$ represents the order of the integrator, and $c_i, d_i \in \bb{R} \ (i = 1,\cdots,k)$ are coefficients chosen such that the residual error is of the order of $(\Delta t)^{k+1}$

In a direct approach, conditions for the coefficients are found by comparing the Taylor expansions of both sides of \eqref{eq:time-evolution_op_prod_approx}
\begin{align*}
	\exp(\Delta t (\hat{D}_K + \hat{D}_U) ) &= \sum_{n=0}^{k} \frac{1}{n!}(\Delta t)^n (\hat{D}_K + \hat{D}_U)^n + \order{(\Delta t)^{k+1}} \\
	\exp(d_i \Delta t \hat{D}_K) &= \sum_{n=0}^{k} \frac{1}{n!}d_i^n(\Delta t)^n \hat{D}_K^n + \order{(\Delta t)^{k+1}} \\
	\exp(c_i \Delta t \hat{D}_U) &= \sum_{n=0}^{k} \frac{1}{n!}c_i^n(\Delta t)^n \hat{D}_U^n + \order{(\Delta t)^{k+1}}
\end{align*}

\subsection{First-Order Integrator ($k=1$)}
\begin{align*}
	\exp(\Delta t (\hat{D}_K + \hat{D}_U) ) &= 1 + \Delta t (\hat{D}_K + \hat{D}_U) + \order{(\Delta t)^2} \\
	\exp(d_i \Delta t \hat{D}_K) &= 1 + d_i \Delta t \hat{D}_K + \order{(\Delta t)^2} \\
	\exp(c_i \Delta t \hat{D}_U) &= 1 + c_i \Delta t \hat{D}_U + \order{(\Delta t)^2} \\
	\exp(d_1 \Delta t \hat{D}_K) \exp(c_1 \Delta t \hat{D}_U) &= 1 + \Delta t (d_1 \hat{D}_K + c_1 \hat{D}_U) + \order{(\Delta t)^2}
\end{align*}

The solution is Lie product formula:
\[ c_1 = d_1 =1 \]
\[ \exp(\Delta t (\hat{D}_K + \hat{D}_U) ) = \exp(\Delta t \hat{D}_K) \exp(\Delta t \hat{D}_U) + \order{(\Delta t)^2} \]
This corresponds to the Symplectic Euler method.

\subsection*{Second-Order Integrator ($k=2$)}
\begin{align*}
	\exp(\Delta t (\hat{D}_K + \hat{D}_U) ) &= 1 + \Delta t (\hat{D}_K + \hat{D}_U) + \frac{1}{2} (\Delta t)^2 (\hat{D}_K^2 + \hat{D}_U \hat{D}_K + \hat{D}_K \hat{D}_U + \hat{D}_U^2) + \order{(\Delta t)^3} \\
	\exp(d_i \Delta t \hat{D}_K) &= 1 + d_i \Delta t \hat{D}_K + \frac{1}{2} d_i^2 (\Delta t)^2 \hat{D}_K^2 + \order{(\Delta t)^3} \\
	\exp(c_i \Delta t \hat{D}_U) &= 1 + c_i \Delta t \hat{D}_U + \frac{1}{2} c_i^2 (\Delta t)^2 \hat{D}_U^2 + \order{(\Delta t)^3} \\
	\prod_{i=2,1} \exp(d_i \Delta t \hat{D}_K) \exp(c_i \Delta t \hat{D}_U) &= 1 + \Delta t \left( (d_1+d_2) \hat{D}_K + (c_1+c_2) \hat{D}_U \right) \\
	&\quad +(\Delta t)^2 \left( \frac{1}{2} (d_1 + d_2)^2 \hat{D}_K^2 + \frac{1}{2} (c_1 + c_2)^2 \hat{D}_U^2 + (d_1 c_1 + d_2 c_2) \hat{D}_K \hat{D}_U + c_2 d_1 \hat{D}_U \hat{D}_K \right)
\end{align*}
\[ d_1 + d_2 = 1 \, , \ c_1 + c_2 = 1 \, , \ d_1 c_1 + d_2 c_2 = \frac{1}{2} \, , \ d_1 c_2 = \frac{1}{2} \]
\begin{gather*}
	d_2 = 1 - d_1 \, , \ c_2 = 1 - c_1 \\
	d_1 c_1 + (1 - d_1)(1 - c_1) = \frac{1}{2} \, , \ d_1 (1 - c_1) = \frac{1}{2} \\
	2d_1 c_1 -d_1 -c_1 = -\frac{1}{2} \, , \ d_1 - \frac{1}{2} = d_1 c_1 \\
	d_1 - c_1 = \frac{1}{2}
\end{gather*}
Two prominent solutions emerge (Strang Splitting):
\[ c_1 = \frac{1}{2} \, , \ d_1 = 1 \, , \ c_2 = \frac{1}{2} \, , \ d_2 = 0 \]
\[ \exp(\Delta t (\hat{D}_K + \hat{D}_U) ) = \exp(\frac{1}{2} \Delta t \hat{D}_U ) \exp(\Delta t \hat{D}_K) \exp(\frac{1}{2} \Delta t \hat{D}_U ) + \order{(\Delta t)^3} \]
which correspons to 'kick-drift-kick' form of leapfrog integration, and

\[ c_1 = 0 \, , \ d_1 = \frac{1}{2} \, , \ c_2 = 1 \, , \ d_2 = \frac{1}{2} \]
\[ \exp(\Delta t (\hat{D}_K + \hat{D}_U) ) = \exp(\frac{1}{2} \Delta t \hat{D}_K ) \exp(\Delta t \hat{D}_U) \exp(\frac{1}{2} \Delta t \hat{D}_K ) + \order{(\Delta t)^3} \]
which correspons to 'drift-kick-drift' form of leapfrog integration.

In general, two conditions emerge from the coefficients of $\hat{D}_K$ and $\hat{D}_U$,
\[ \sum_{i=1}^{k} c_i = 1 \, , \ \sum_{i=1}^{k} d_i = 1 \]
and higher order terms provide additional conditions.

However, as noted by Yoshida: \textit{``with this direct method, it is almost hopeless to obtain a much higher integrator''} \cite{YOSHIDA1990262}.
Instead, coefficient sets for higher-order integrators (e.g., $k=4$) are systematically derived using the Baker--Campbell--Hausdorff formula (see Yoshida \cite{YOSHIDA1990262}).

\subsection{Action on Phase-Space Vectors}
Because the Hamiltonian is separable,
\begin{gather*}
	\dot{q}^i = \pdv{H}{p_i} = \pdv{K}{p_i} \\
	\dot{p}_i = -\pdv{H}{q^i} = -\pdv{U}{q^i}
\end{gather*}
\begin{align*}
	\hat{D}_K \vec{z} = \{ \vec{z}, K \} =
	\begin{pmatrix}
		\{ \vec{q}, K \} \\
		\{ \vec{p}, K \}
	\end{pmatrix} 
\end{align*}
\begin{align*}
	\{ \vec{q}, K \} &= 
	\begin{pmatrix}
		\{ q_1, K \} \\
		\vdots \\
		\{ q_n, K \}
	\end{pmatrix} \\
	&=
	\begin{pmatrix}
		\sum_i \left( \pdv{q_1}{q^i} \pdv{K}{p_i} - \pdv{K}{q^i} \pdv{q_1}{p_i} \right) \\
		\vdots \\
		\sum_i \left( \pdv{q_n}{q^i} \pdv{K}{p_i} - \pdv{K}{q^i} \pdv{q_n}{p_i} \right)
	\end{pmatrix} \\
	&=
	\begin{pmatrix}
		\pdv{K}{p_1} \\
		\vdots \\
		\pdv{K}{p_n}
	\end{pmatrix} =
	\begin{pmatrix}
		\dot{q}^1 \\
		\vdots \\
		\dot{q}^n
	\end{pmatrix}
	= \dot{\vec{q}}
\end{align*}
\begin{align*}
	\{ \vec{p}, K \} &= 
	\begin{pmatrix}
		\{ p_1, K \} \\
		\vdots \\
		\{ p_n, K \}
	\end{pmatrix} \\
	&=
	\begin{pmatrix}
		\sum_i \left( \pdv{p_1}{q^i} \pdv{K}{p_i} - \pdv{K}{q^i} \pdv{p_1}{p_i} \right) \\
		\vdots \\
		\sum_i \left( \pdv{p_n}{q^i} \pdv{K}{p_i} - \pdv{K}{q^i} \pdv{p_n}{p_i} \right)
	\end{pmatrix} \\
	&=
	\begin{pmatrix}
		0 \\
		\vdots \\
		0
	\end{pmatrix} = \vec{0}
\end{align*}
\begin{equation}
	\hat{D}_K \vec{z} = \{ \vec{z}, K \} =
\begin{pmatrix}
	\dot{\vec{q}} \\
	\vec{0}
\end{pmatrix}
\end{equation}
\begin{equation}
	\hat{D}_K^2 \vec{z} = \{ \{ \vec{z}, K \}, K \} =
	\begin{pmatrix}
		\{ \dot{\vec{q}}, K \} \\
		\{ \vec{0}, K \}
	\end{pmatrix} = \vec{0}
\end{equation}

\begin{align*}
	\hat{D}_U \vec{z} = \{ \vec{z}, U \} =
	\begin{pmatrix}
		\{ \vec{q}, U \} \\
		\{ \vec{p}, U \}
	\end{pmatrix} 
\end{align*}
\begin{align*}
	\{ \vec{q}, U \} &= 
	\begin{pmatrix}
		\{ q_1, U \} \\
		\vdots \\
		\{ q_n, U \}
	\end{pmatrix} \\
	&=
	\begin{pmatrix}
		\sum_i \left( \pdv{q_1}{q^i} \pdv{U}{p_i} - \pdv{U}{q^i} \pdv{q_1}{p_i} \right) \\
		\vdots \\
		\sum_i \left( \pdv{q_n}{q^i} \pdv{U}{p_i} - \pdv{U}{q^i} \pdv{q_n}{p_i} \right)
	\end{pmatrix} \\
	&=
	\begin{pmatrix}
		0 \\
		\vdots \\
		0
	\end{pmatrix} = \vec{0}
\end{align*}
\begin{align*}
	\{ \vec{p}, U \} &= 
	\begin{pmatrix}
		\{ p_1, U \} \\
		\vdots \\
		\{ p_n, U \}
	\end{pmatrix} \\
	&=
	\begin{pmatrix}
		\sum_i \left( \pdv{p_1}{q^i} \pdv{U}{p_i} - \pdv{U}{q^i} \pdv{p_1}{p_i} \right) \\
		\vdots \\
		\sum_i \left( \pdv{p_n}{q^i} \pdv{U}{p_i} - \pdv{U}{q^i} \pdv{p_n}{p_i} \right)
	\end{pmatrix} \\
	&=
	\begin{pmatrix}
		-\pdv{U}{q^1} \\
		\vdots \\
		-\pdv{U}{q^n}
	\end{pmatrix} =
	\begin{pmatrix}
		\dot{p}_1 \\
		\vdots \\
		\dot{p}_n
	\end{pmatrix}
	= \dot{\vec{p}}
\end{align*}
\begin{equation}
	\hat{D}_U \vec{z} = \{ \vec{z}, U \} =
\begin{pmatrix}
	\vec{0} \\
	\dot{\vec{p}}
\end{pmatrix}
\end{equation}
\begin{equation}
	\hat{D}_U^2 \vec{z} = \{ \{ \vec{z}, U \}, U \} =
	\begin{pmatrix}
		\{ \vec{0}, U \} \\
		\{ \dot{\vec{p}}, U \}
	\end{pmatrix} = \vec{0}
\end{equation}

Using a Taylor expansion, for $\alpha \in \bb{R}$,
\begin{align*}
	\exp(\alpha\hat{D}_U) \vec{z} &= \sum_{n=0}^{\infty} \frac{\alpha^n\hat{D}_U^n}{n!} \vec{z} = (1 + \alpha\hat{D}_U) \vec{z} \\
	\exp(\alpha\hat{D}_K) \vec{z} &= \sum_{n=0}^{\infty} \frac{\alpha^n\hat{D}_K^n}{n!} \vec{z} = (1 + \alpha\hat{D}_K) \vec{z}
\end{align*}

\begin{align*}
	\exp(c_i \Delta t \hat{D}_U) \begin{pmatrix} \vec{q} \\ \vec{p} \end{pmatrix} &= \begin{pmatrix} \vec{q} \\ \vec{p} + c_i \Delta t \dot{\vec{p}}\end{pmatrix} \\
	\exp(d_i \Delta t \hat{D}_K) \begin{pmatrix} \vec{q} \\ \vec{p} \end{pmatrix} &= \begin{pmatrix} \vec{q} + d_i \Delta t \dot{\vec{q}} \\ \vec{p} \end{pmatrix}
\end{align*}

\bibliographystyle{plain}
\bibliography{refs}

\end{document}